% -------------------------------------------------------------
% thesis_style.tex
% This file defines custom style updates for the PhD thesis,
% including colors and fonts.
% -------------------------------------------------------------

% 1. Define custom colors for the document
% You can change these RGB values to match your university's branding
\definecolor{ThesisDarkRed}{RGB}{140, 0, 0}
\definecolor{ThesisAccentRed}{RGB}{210, 0, 0}
\definecolor{ThesisGray}{RGB}{100, 100, 100}

% 2. Update font formatting for headings using KOMA-Script commands
% This applies the dark blue color to chapters, sections, etc.
\addtokomafont{chapter}{\color{ThesisDarkRed}}
\addtokomafont{section}{\color{ThesisDarkRed}}
\addtokomafont{subsection}{\color{ThesisAccentRed}}
\addtokomafont{subsubsection}{\color{ThesisAccentRed}}
\addtokomafont{paragraph}{\color{ThesisAccentRed}}
\addtokomafont{captionlabel}{\color{ThesisDarkRed}\bfseries}
\addtokomafont{pagenumber}{\color{ThesisGray}}

% 3. Update hyperlink colors
% Here we apply the custom colors to links, citations, and URLs
\hypersetup{
    colorlinks=true,
    linkcolor=black,
    filecolor=black,
    urlcolor=ThesisDarkRed,
    citecolor=ThesisDarkRed,
}

% 4. Enable chapter formatting
% kaobook supports different chapter styles (e.g. kao, plain)
\setchapterstyle{kao}

% 5. Cross-referencing helpers using cleveref
% This ensures consistent formatting (e.g., "Figure" vs "Fig." etc.) across the document.
% By default kaobook uses noabbrev for cleveref. If you prefer abbreviations
% (so \cref outputs Fig. instead of Figure), uncomment the next lines:
% \crefname{figure}{Fig.}{Figs.}
% \Crefname{figure}{Fig.}{Figs.}
% \crefname{equation}{Eq.}{Eqs.}
% \Crefname{equation}{Eq.}{Eqs.}
% \crefname{table}{Tab.}{Tabs.}
% \Crefname{table}{Tab.}{Tabs.}
% \crefname{section}{Sec.}{Secs.}
% \Crefname{section}{Sec.}{Secs.}

% Standardized helper macros bridging to \cref
\newcommand{\figref}[1]{\cref{#1}}
\newcommand{\tabref}[1]{\cref{#1}}
\newcommand{\secref}[1]{\cref{#1}}
\newcommand{\chapref}[1]{\cref{#1}}
\newcommand{\eqnref}[1]{\cref{#1}}
\newcommand{\appref}[1]{\cref{#1}}
